\section*{Conclusion}

In this work, we take one of the first looks at the genomic signatures of evolution
in large populations using the Pacific acorn barnacle, \balgla{}.
In accordance with population genetics theory, we find evidence for the increased
efficacy of selection in this large barnacle population, manifested as both purifying
selection reducing diversity and maintaining conservation across functional sites, and
pervasive positive selection driving adaptive divergence in protein-coding genes.
Surprisingly, however, we find a highly repetitive genome, in apparent contrast to
the expectation that purifying selection should limit the proliferation of repetitive
sequences.
\comment{From Scott: not sure right now is the best place to bring this up for the first
time. Is it the first time?}
At the population scale, we observe average genome-wide
nucleotide diversity exceeding 5\%, placing this species among hyperpolymorphic
organisms.
Yet, while this diversity implies large effective population sizes ($\approx10^{6}$
individuals), it still pales in comparison to the estimated census sizes
of this population, underscoring the long-standing disparity between effective
and census population sizes.
Together, these findings establish \bgla{} as a promising emerging model for studying
the mechanisms of evolution at biological extremes.

